\documentclass[11pt]{article}
\usepackage[a4paper,margin=2.4cm]{geometry}
\usepackage[T1]{fontenc}
\usepackage{lmodern}
\usepackage{amsmath,amssymb}
\usepackage{booktabs}
\usepackage{graphicx}
\usepackage{xcolor}
\usepackage{hyperref}
\usepackage[numbers,sort&compress]{natbib}
\usepackage{caption}
\captionsetup{font=small,labelfont=bf}
\newcommand{\todo}[1]{\textcolor{red}{[TODO: #1]}}
\newcommand{\GEARS}{\textsc{GEARS}}
\newcommand{\LASSER}{\textsc{LASSER}}

\title{Does the gene co-expression graph matter in GEARS?\\
\large Learned, density-matched and random graphs give the same perturbation predictions}

\author{\todo{author names and affiliations}}
\date{}

\begin{document}
\maketitle

\begin{abstract}
Graph neural networks for predicting transcriptional responses to genetic perturbations rely on gene--gene graphs as
priors. \GEARS{}~\cite{roohani2024gears} combines a Gene Ontology graph over perturbed genes with a gene
co-expression graph over output genes. A natural hypothesis is that a better co-expression graph should improve
prediction. We test it with \LASSER{}, a self-supervised, multi-view graph-structure-learning method in the spirit
of GSR~\cite{zhao2023gsr}, under a leakage-safe protocol with exact parity to upstream \GEARS{}.
We report three findings on the Norman et al.\ combinatorial CRISPRa screen~\cite{norman2019}.
\textbf{(i)}~The co-expression module of the reference implementation is largely inert. Its Pearson graph leaves
87\% of genes without neighbours, and a batching defect means that only the first cell of each mini-batch receives
graph messages.
\textbf{(ii)}~Once the defect is fixed, the learned graph, a random graph with identical degrees and weights, and a
density-matched Pearson graph all perform like the original sparse graph. Learned minus original Pearson~$\Delta$:
$-0.013$, 95\% CI $[-0.052, 0.027]$, 0/3 splits better. All differences are smaller than the run-to-run
variability of a single configuration ($\approx 0.08$).
\textbf{(iii)}~This is not because the learned graph is uninformative. Its new edges are co-expressed in held-out
cells 18$\times$ more often than random pairs, and the edges that recur across splits form coherent modules
(mitosis, proteasome, nucleolar/ribosome biogenesis).
We trace the null result to the architecture. In \GEARS{} the co-expression GNN smooths a learnable embedding table
that does not depend on the cell's expression, so any graph acts as a re-parameterisation of free per-gene
parameters. We argue that graph priors over output genes can only matter if they propagate cell- or
perturbation-specific signals, and we give recommendations for evaluating graph priors in perturbation models.
\end{abstract}

% =====================================================================================
\section{Introduction}

Predicting the transcriptional outcome of unseen genetic perturbations, in particular combinations, is a central
task for in-silico screening. \GEARS{}~\cite{roohani2024gears} is a widely used graph-based approach. It encodes
prior knowledge through two graphs. A Gene Ontology (GO) similarity graph over perturbable genes lets unseen
perturbations borrow from functionally related seen ones. A co-expression graph over output genes, built once from
Pearson correlations in training cells, is meant to share information between genes that move together.

The co-expression graph is an obvious target for improvement. It is fixed, built by thresholding noisy
single-cell correlations, and blind to the task. Graph structure learning (GSL) methods, either jointly with the
task (e.g.\ IDGL~\cite{chen2020idgl}, LDS~\cite{franceschi2019lds}) or self-supervised and decoupled from it
(GSR~\cite{zhao2023gsr}), have improved node-level prediction elsewhere. In preliminary work we tried task-driven
GSL (IDGL-style learning of a cosine-similarity graph from embeddings) and hyperbolic GO embeddings
(Poincar\'e~\cite{nickel2017poincare}) inside \GEARS{}, and saw no gains
\todo{one or two sentences with the exact settings and numbers of these earlier attempts}. Those attempts also had
practical problems: dense learned graphs that didn't fit in memory, unstable optimisation and embedding collapse.
That motivated the decoupled, memory-safe design studied here.

This paper asks a sharper question: \emph{does the co-expression graph matter at all in \GEARS{}, and if not, why?}
Our contributions are:
\begin{enumerate}
  \item A leakage-safe experimental harness around \GEARS{} that reproduces the upstream model bit-for-bit when graph
        learning is off. Parity is verified on graph tensors, initial weights, the first training losses and the
        final model.
  \item \LASSER{}, a self-supervised multi-view graph-structure learner for output genes (pretrain, refine once,
        then train \GEARS{} unchanged). It uses only training cells and is memory-safe (no dense $N{\times}N$ tensor
        in autograd).
  \item The identification of a batching defect in the reference implementation that confines co-expression
        message passing to the first cell of each mini-batch, and an opt-in fix.
  \item A controlled comparison, with the batching fix, of the original Pearson graph, the learned graph, a
        density-matched Pearson graph and a degree- and weight-preserving random rewiring. All perform alike, within
        run-to-run variability.
  \item A mechanistic explanation. In \GEARS{} the co-expression GNN never sees cell-specific input, so graph
        structure can only re-parameterise free per-gene embeddings. That explains why no graph-learning paradigm
        we tried changed the predictions.
\end{enumerate}

% =====================================================================================
\section{Background}

\paragraph{GEARS.} For a control cell with expression $x\in\mathbb{R}^{N}$ ($N{=}5{,}045$ genes) and a set of
perturbed genes, \GEARS{} computes per-gene embeddings from two learnable tables:
\begin{equation}
  h_g = \mathrm{MLP}\!\left(\mathrm{ReLU}(\mathrm{BN}(E_g)) + 0.2\,\big[\mathrm{SGConv}_{A}(P)\big]_g\right),
  \label{eq:gears}
\end{equation}
where $E, P \in \mathbb{R}^{N\times d}$ are embedding tables and $A$ is the co-expression graph. Perturbation
embeddings, propagated over the GO graph, are added to every gene of the perturbed cell. Gene-specific decoders and
a cross-gene MLP then produce the predicted change, and $x$ is added as a residual at the end. Two properties of the
reference implementation are central to this paper:
\begin{itemize}
  \item \textbf{The co-expression path is cell-independent.} The input to $\mathrm{SGConv}_A$ is the table $P$, the
        same for every cell. The cell's expression $x$ enters only through the final residual. So
        $\mathrm{SGConv}_A(P)$ is a fixed matrix of per-gene vectors, a smooth function of free parameters.
  \item \textbf{Batching defect.} Embeddings of the $B$ cells in a mini-batch are stacked into $BN$ nodes, but
        $\mathrm{SGConv}_A$ is called with the single-cell \texttt{edge\_index} (node ids $0,\dots,N{-}1$). Only the
        first cell's copy aggregates neighbours; the other $B{-}1$ copies receive only their own self-loop, i.e.\ a
        linear map of $P$.
\end{itemize}

\paragraph{Graph structure learning.} Task-driven GSL learns the graph jointly with the downstream
loss~\cite{chen2020idgl,franceschi2019lds}. GSR~\cite{zhao2023gsr} instead pretrains multi-view node encoders by
contrastive link prediction~\cite{oord2018infonce,he2020moco}, refines the graph once by adding and removing edges
by embedding similarity, and then trains the task model on the fixed refined graph. This decoupling is attractive
here because the task offers few independent training signals (138 training perturbations per split).

\paragraph{Baselines in perturbation prediction.} Recent benchmarks report that deep perturbation models,
\GEARS{} included, often do not outperform simple additive or linear baselines~\cite{ahlmanneltze2025linear,
csendes2024benchmark}. Our results add a mechanistic reason why one of \GEARS{}' two priors contributes little.

% =====================================================================================
\section{Methods}

\subsection{Data, splits and evaluation}
We use the Norman et al.\ CRISPRa screen in K562 cells~\cite{norman2019} as processed by \GEARS{}: 5,045 genes,
single and double perturbations. We use \GEARS{}' \emph{simulation} split (seeds 1--5,
\texttt{train\_gene\_set\_size}=0.75). Test perturbations are grouped into unseen singles and combinations with
0, 1 or 2 genes seen in training. In split~1, for example, there are 36, 9, 43 and 19 such perturbations, and 138
training perturbations. Following the \GEARS{} paper, we compute metrics on mean predicted and mean observed
expression per test perturbation:
\begin{itemize}
  \item Pearson correlation of predicted vs.\ observed change from control, over all genes ($\text{Pearson}\,\Delta$)
        and over the top-20 differentially expressed (DE) genes;
  \item MSE on the top-20 DE genes, and its value normalised by the no-perturbation baseline (ratio of means);
  \item the fraction of top-20 DE genes predicted with the wrong sign.
\end{itemize}
Reference baselines are scored with the same code: \emph{no-perturbation} (predict control) and \emph{additive}
(control plus the sum of the training single-gene changes). The additive baseline equals no-perturbation for unseen
singles. Note that both baselines predict exactly zero change for some genes, which counts as neither correct nor
opposite direction, so their low opposite-direction fractions are not comparable with the models'.

\subsection{Harness and leakage control}
All arms train the unmodified \GEARS{} model, loss, training loop and default hyperparameters (hidden size 64,
Adam with learning rate $10^{-3}$, weight decay $5{\times}10^{-4}$, batch size 32). The epoch with the lowest
validation top-20 DE MSE is kept. Everything is implemented in PyTorch and PyTorch Geometric~\cite{fey2019pyg}, with \GEARS{} vendored unchanged.
The only intervention is the co-expression graph passed to
\texttt{model\_initialize} (plus, in the second study, the opt-in batching fix). With graph learning off, a test
suite verifies exact equality with upstream \GEARS{} in the co-expression and GO graph tensors, the initial
\texttt{state\_dict}, the first training losses and the best model after an epoch. Everything the graph learner
sees (features, positives and the base graph) comes from control cells and training-perturbation cells of the
current split only. A test that permutes and rescales validation and test cells leaves the learned graph
bit-for-bit unchanged.

\subsection{LASSER: self-supervised co-expression graph refinement}
\paragraph{Base graph $A_0$.} The \GEARS{} Pearson graph: for each gene, up to 20 neighbours with $|r|>0.4$,
computed on control cells plus training single-perturbation cells.

\paragraph{Views.} For each gene we build three standardised feature vectors from training data only:
\begin{itemize}
  \item E: expression over 500 $k$-means metacells, reduced by PCA to 64 dimensions;
  \item R: mean change from control across the training perturbations, PCA to 32 dimensions;
  \item S: DeepWalk~\cite{perozzi2014deepwalk} embedding of $A_0$, 64 dimensions.
\end{itemize}
Genes with no signal in a view (zero variance, or isolated in $A_0$) get a learned \emph{missing} token instead of
a zero vector.

\paragraph{Pretraining.} The training cells are split into halves $H_1/H_2$, stratified by condition.
\begin{itemize}
  \item Positives: Pearson top-20 pairs on $H_1$ that stay above the threshold in at least 70\% of 20 bootstrap
        resamples.
  \item Negatives: 256 per query, a quarter of them hard (high embedding similarity but low $|r|$ on $H_1$),
        excluding $A_0$ and positive neighbours.
  \item Encoders: per view, an MLP, two SGConv layers~\cite{wu2019sgc}, then a linear layer with LayerNorm (no
        output ReLU). Keys come from momentum encoders ($m{=}0.99$).
  \item Loss: intra- and inter-view InfoNCE ($\tau{=}0.2$, intra weight $\alpha{=}0.75$) plus a VICReg variance
        term~\cite{bardes2022vicreg}.
  \item Training: each epoch, positives are split 70/30 into message edges and target edges, so a target edge is
        never in the GNN input. Early stopping is on held-out link AUC.
\end{itemize}

\paragraph{Refinement (once).} Edge scores are a weighted sum ($\beta = 0.4, 0.3, 0.3$ for E, R, S) of min--max
normalised cosine similarities, computed in chunks without gradients. Per gene, we remove the 10\% weakest $A_0$
edges and add the top 5 new candidates, top up every gene to in-degree $\ge 5$, and cap new in-edges at 25. Kept
edges keep their Pearson weight; new edges get $0.5\times$ their score quantile-mapped onto the $A_0$ weight
distribution. On a T4 GPU the whole procedure takes about 25\,s and under 1.6\,GB.

\subsection{Controls and the batching fix}
\begin{itemize}
  \item \textbf{GEARS-kNN8}: the \GEARS{} Pearson graph without threshold, 8 neighbours per gene, i.e.\ about the
        density of the learned graph.
  \item \textbf{LASSER-random}: the learned graph with every new edge rewired to a uniformly random source gene
        (no self-loops, $A_0$ pairs or duplicates). This keeps each gene's in-degree and the weight distribution but
        destroys \emph{which} genes are connected.
  \item \textbf{Batching fix} (opt-in, ``-fix''): the co-expression SGConv repeats the graph once per cell with
        offset node ids. Parameters are unchanged, and with one cell per batch the output is identical to upstream.
\end{itemize}

\subsection{Studies and statistics}
\textbf{Study 1 (upstream \GEARS{}, 20 epochs)}: \GEARS{} on splits 2--5 and \LASSER{} on splits 1--3.
\textbf{Study 2 (batching fix, 10 epochs)}: \GEARS-fix, \LASSER-fix and \LASSER-random-fix on splits 1--3, and
\GEARS-kNN8-fix on split 1. Within a split, all arms share the model initialisation and batch order, so
differences isolate the graph.
\begin{itemize}
  \item Across splits: mean $\pm$ sd and the paired improvement with a $t$-based 95\% CI. With $n{=}3$ splits,
        $p$-values carry little information, so we report effect sizes and the number of splits improved.
  \item Per perturbation: (split, perturbation) pairs, the share improved, Wilcoxon signed-rank, and a 95\%
        bootstrap CI that resamples splits and then perturbations within splits.
\end{itemize}
Hardware: one NVIDIA T4 per run (Kaggle). A 10-epoch run takes 42--50 minutes.

% =====================================================================================
\section{Results}

\subsection{The co-expression prior of GEARS is sparse and mostly unused}
Across splits, the \GEARS{} Pearson graph has 5,992--6,364 edges, of which 4,685--4,726 are self-loops.
4,419--4,495 of the 5,045 genes (87--89\%) have no neighbour at all (Fig.~\ref{fig:graphs}a), because
$|r|>0.4$ is rare in sparse single-cell data. Combined with the batching defect, this means that in the reference
implementation co-expression messages reach about 12\% of genes, and only for 1 of 32 cells per mini-batch.

\subsection{The learned graph is informative}
\LASSER{} produces about 45k edges per split (44,610--45,192). Every gene gets in-degree $\ge 5$ (median 7, max
25), and neighbourhood Jaccard overlap with $A_0$ is 0.03.
\begin{itemize}
  \item \emph{Link prediction}: the pretrained encoders reach a held-out link AUC of 0.973--0.997, vs.\
        0.956--0.979 for the raw view features.
  \item \emph{Held-out co-expression}: 3.4--3.7\% of new edges have $|r|>0.2$ on the held-out half $H_2$, vs.\
        0.1--0.2\% for degree-matched random pairs (about 18-fold enrichment) and 58--66\% for $A_0$ edges
        (Fig.~\ref{fig:graphs}b). The new edges are weaker than $A_0$'s, as intended, since most connect previously
        isolated genes.
  \item \emph{Stability}: 346 new pairs appear in all three splits, although the training perturbations differ
        between splits. They form biologically coherent modules (Table~\ref{tab:consensus}), including mitosis and
        chromosome segregation (\textit{CDCA8}, \textit{CENPA}, \textit{AURKA}, \textit{PLK1}, \textit{TPX2},
        \textit{CDK1}, \textit{NUSAP1}, \textit{MKI67}, \textit{BIRC5}, \textit{CKS2}), the proteasome
        (\textit{PSMA4}, \textit{PSMB1--3}, \textit{UBB}), nucleolar and ribosome biogenesis (\textit{NCL},
        \textit{DDX21}, \textit{SRM}, \textit{HSPD1}) and haematopoietic adhesion (\textit{ITGB2}, \textit{CD84},
        \textit{MYO1F}).
  \item \emph{Beyond Pearson}: 44.5\% of the consensus pairs are also in the density-matched Pearson graph of
        split 1, so the learner recovers correlation structure and adds pairs that pairwise correlation ranks
        lower.
\end{itemize}

\subsection{Study 1: with upstream GEARS, the graph does not change predictions}
On the splits where both arms were run, \GEARS{} and \LASSER{} give Pearson~$\Delta$ of 0.524 vs.\ 0.523
(split~2) and 0.547 vs.\ 0.551 (split~3). Per perturbation (226 pairs), the mean improvement is $+0.002$, 95\% CI
$[-0.003, 0.006]$, with 54\% of perturbations improved. This is what the batching defect predicts.

\subsection{Study 2: with the fix, graph structure still does not matter}
Table~\ref{tab:main} and Fig.~\ref{fig:headline} summarise the batching-fix study. All graph variants are well above
the additive baseline in Pearson~$\Delta$ (0.53 vs.\ 0.41) and in top-20 DE MSE (0.22 vs.\ 0.27), and they are
indistinguishable from each other:
\begin{itemize}
  \item \textbf{Learned vs.\ original Pearson.} \LASSER-fix minus \GEARS-fix: Pearson~$\Delta$ $-0.013$
        $[-0.052, 0.027]$, 0/3 splits better; top-20 DE MSE $-0.002$ $[-0.027, 0.023]$ (positive = better).
        Per perturbation (333 pairs): $-0.012$, bootstrap 95\% CI $[-0.029, 0.001]$, 43\% improved. If anything,
        the learned graph is slightly worse, and the effect is small (Fig.~\ref{fig:perpert}).
  \item \textbf{Learned vs.\ random partners.} \LASSER-fix minus \LASSER-random-fix: Pearson~$\Delta$ $+0.002$
        $[-0.081, 0.084]$; per perturbation $+0.001$ $[-0.028, 0.035]$. Keeping the degrees and weights but
        destroying \emph{which} genes are connected changes nothing measurable.
  \item \textbf{Density-matched Pearson} (split 1 only): \LASSER-fix $0.535$ vs.\ \GEARS-kNN8-fix $0.509$, with
        near-identical MSE.
  \item \textbf{Subgroups} (Fig.~\ref{fig:subgroups}): no subgroup shows a consistent gain. Combinations with both
        genes seen have the largest point estimate in favour of learned over random partners (+0.019 Pearson~$\Delta$
        per perturbation, 65\% improved), but its CI includes zero.
  \item \textbf{Gene groups} (Fig.~\ref{fig:genegroups}): if a denser graph helped, the 4.4k genes isolated in
        $A_0$, whose neighbourhoods change most, should gain. They do not ($-0.015 \pm 0.021$ Pearson~$\Delta$ over
        splits).
\end{itemize}

\subsection{Run-to-run variability exceeds the graph effect}
Two runs of the same \LASSER{} configuration on split~1, in separate sessions, gave Pearson~$\Delta$ of 0.574 and
0.491. The graphs differed by 0.2\% of edges because of a numerical change, and GPU kernels were not
deterministic. Within study 2, the spread across all graph variants on a split was 0.067, 0.002 and 0.030 for
splits 1--3. Single-split comparisons of graph priors in \GEARS{} are therefore unreliable; the +0.04 gain we saw
in an early single-split run did not replicate.

% =====================================================================================
\section{Why the graph cannot matter much in GEARS}
Equation~\ref{eq:gears} makes the reason explicit. With the batching fix, the co-expression path contributes
$Q = \mathrm{SGConv}_A(P) = \hat{A} P \Theta + b$ to every cell, where $P$ is a free learnable table and
$\hat{A}$ is the normalised adjacency with self-loops. $Q$ depends neither on the cell's expression nor on the
perturbation. Whenever $\hat{A}$ is invertible (with self-loops this is the generic case), any target $Q$ is
reachable by choosing $P = \hat{A}^{-1}(Q-b)\Theta^{+}$, for every graph. The graph therefore changes only the
parameterisation of a per-gene bias, i.e.\ the optimisation path and the implicit regularisation (max-norm
embeddings, weight decay), not what the model can express. Our controls are consistent with this: original, dense,
learned and random graphs end up equivalent. Without the fix, the first cell of each batch gets
$\hat{A}P\Theta+b$ and the others $P\Theta+b$, which is still cell-independent.

More broadly, the predicted change in \GEARS{} is computed from the perturbation and learnable per-gene parameters.
The control cell's expression is added at the end as a residual, apart from batch-norm statistics during training.
Gene--gene priors over output genes can only add information if messages carry cell- or perturbation-specific
signals: for example, propagating the perturbation's effect or the cell's state over the graph. This also explains
why task-driven graph learning (IDGL-style) and hyperbolic GO embeddings, which we tried earlier, did not help
either.

% =====================================================================================
\section{Recommendations}
\begin{enumerate}
  \item \textbf{Check that the prior can act.} Before learning a graph, verify that it reaches the computation
        (here, the batching defect) and that its input is not a free, cell-independent table (here, the
        re-parameterisation argument).
  \item \textbf{Always include structure-destroying controls.} A degree- and weight-preserving random rewiring and a
        density-matched simple graph separate ``learned structure helps'' from ``more edges help'' and from noise.
  \item \textbf{Report variability and pair the comparisons.} Use several splits and training seeds, pair arms
        within split and seed, and report effect sizes with CIs and per-perturbation analyses rather than single
        runs.
  \item \textbf{Report the reference baselines.} Include the additive and no-perturbation baselines scored with the
        same code, and interpret direction metrics with care for baselines that predict zero change.
\end{enumerate}

% =====================================================================================
\section{Limitations}
\begin{itemize}
  \item One dataset (Norman K562).
  \item Three splits in the main study, 10 epochs per run (chosen to fit the compute budget; validation DE-MSE
        usually peaks early), and only one split for the density-matched control.
  \item Run-to-run variability was estimated across sessions rather than from planned training-seed replicates.
  \item Planned ablations (single views, add-only refinement, embedding initialisation from pretrained views,
        edit sizes) were not run. Given the re-parameterisation argument, we do not expect them to change the
        conclusion for \GEARS{}, but they would matter for architectures where the graph acts on cell-specific
        signals.
  \item External validation of consensus edges (e.g.\ STRING, TRRUST) is left to future work.
\end{itemize}

\section{Conclusion}
In \GEARS{}, the gene co-expression graph is a nearly empty prior, reaches only one cell per mini-batch, and, even
when fixed, cannot change what the model can express. A self-supervised learned graph captures real and stable
co-expression structure, but neither it, nor a density-matched Pearson graph, nor a random rewiring changes
predictions beyond run-to-run noise. Improving graph priors for perturbation prediction will require architectures
in which graph messages carry cell- or perturbation-specific information, evaluated against structure-destroying
controls.

\paragraph{Code and data availability.} Code, tests, run configurations and per-run outputs:
\todo{repository URL}. The harness reproduces upstream \GEARS{} exactly when graph learning and the fix are off.

% =====================================================================================
\begin{table}[t]
\centering
\caption{Study 2 (batching fix, 10 epochs): test metrics, mean $\pm$ sd over splits 1--3 (GEARS-kNN8-fix: split 1
only). Pearson~$\Delta$: higher is better; MSE, nMSE and opposite direction: lower is better. nMSE = top-20 DE MSE
divided by the no-perturbation MSE (ratio of means). $^\dagger$Baselines that predict zero change for some genes
have a trivially low opposite-direction fraction.}
\label{tab:main}
\small
\begin{tabular}{lccccc}
\toprule
 & Pearson $\Delta$ & Pearson $\Delta$ (top-20 DE) & MSE (top-20 DE) & nMSE & Opposite dir.\\
\midrule
No-perturbation & 0.000 & 0.000 & $0.469\pm0.043$ & 1.000 & $0.000^\dagger$\\
Additive & $0.407\pm0.024$ & $0.448\pm0.022$ & $0.268\pm0.033$ & $0.570\pm0.025$ & $0.083\pm0.010^\dagger$\\
\midrule
GEARS-fix (Pearson $A_0$) & $0.540\pm0.023$ & $0.670\pm0.017$ & $0.218\pm0.040$ & $0.462\pm0.042$ & $0.350\pm0.011$\\
LASSER-fix (learned) & $0.527\pm0.008$ & $0.671\pm0.023$ & $0.220\pm0.032$ & $0.467\pm0.025$ & $0.347\pm0.011$\\
LASSER-random-fix & $0.526\pm0.029$ & $0.657\pm0.017$ & $0.221\pm0.034$ & $0.471\pm0.034$ & $0.346\pm0.006$\\
GEARS-kNN8-fix (1 split) & 0.509 & 0.672 & 0.192 & 0.439 & 0.369\\
\bottomrule
\end{tabular}
\end{table}

\begin{table}[t]
\centering
\caption{Paired differences (arm minus reference, oriented so that $>0$ is better), study 2.
Splits: mean over splits with a $t$-based 95\% CI and the number of splits improved.
Perturbations: mean over (split, perturbation) pairs with a 95\% bootstrap CI that resamples splits.}
\label{tab:paired}
\small
\begin{tabular}{llccc}
\toprule
Comparison & Metric & Over splits & Splits better & Over perturbations\\
\midrule
LASSER-fix vs GEARS-fix & Pearson $\Delta$ & $-0.013\;[-0.052, 0.027]$ & 0/3 & $-0.012\;[-0.029, 0.001]$\\
 & MSE (top-20 DE) & $-0.002\;[-0.027, 0.023]$ & 1/3 & $-0.002\;[-0.012, 0.009]$\\
 & Opposite direction & $+0.002\;[-0.019, 0.024]$ & 2/3 & $+0.002\;[-0.008, 0.013]$\\
LASSER-fix vs LASSER-random-fix & Pearson $\Delta$ & $+0.002\;[-0.081, 0.084]$ & 1/3 & $+0.001\;[-0.028, 0.035]$\\
 & MSE (top-20 DE) & $+0.002\;[-0.034, 0.037]$ & 2/3 & $+0.002\;[-0.013, 0.014]$\\
 & Opposite direction & $-0.002\;[-0.013, 0.009]$ & 1/3 & $-0.002\;[-0.010, 0.006]$\\
LASSER-fix vs GEARS-kNN8-fix & Pearson $\Delta$ & $+0.026$ (1 split) & 1/1 & $+0.026\;[0.013, 0.041]$\\
\bottomrule
\end{tabular}
\end{table}

\begin{table}[t]
\centering
\caption{The 12 highest-scoring learned edges present in all three splits (346 such pairs in total; full list in
\texttt{consensus\_edges\_split1-3.csv}). ``kNN8'': the pair is also in the density-matched Pearson graph of split 1.}
\label{tab:consensus}
\small
\begin{tabular}{llcl}
\toprule
Gene & Gene & kNN8 & Module\\
\midrule
CDCA8 & CENPA & yes & mitosis / chromosome segregation\\
CENPA & AURKA & yes & mitosis\\
CKS2 & BIRC5 & yes & mitosis / cell cycle\\
PLK1 & TPX2 & yes & mitotic spindle\\
PSMB2 & PSMB3 & yes & proteasome\\
SRM & NCL & yes & nucleolus / growth\\
PSMB2 & PSMA4 & yes & proteasome\\
MYO1F & ITGB2 & no & haematopoietic adhesion\\
ENO1 & HSPD1 & yes & metabolism / chaperone\\
HSPD1 & DDX21 & yes & nucleolus / ribosome biogenesis\\
CDK1 & NUSAP1 & yes & mitosis\\
CD84 & ITGB2 & no & haematopoietic adhesion\\
\bottomrule
\end{tabular}
\end{table}

\begin{figure}[t]
\centering
\includegraphics[width=\linewidth]{figures/fig1_headline_overall.pdf}
\caption{Study 2: test metrics per arm. Dots: splits; marker: mean $\pm$ 95\% CI; grey lines join arms within a
split. All graph variants are similar and clearly above the reference baselines in Pearson~$\Delta$ and MSE.}
\label{fig:headline}
\end{figure}

\begin{figure}[t]
\centering
\includegraphics[width=\linewidth]{figures/fig_forest_fix.pdf}
\caption{Paired improvement over splits (oriented so that $>0$ is better), with 95\% CI and the number of splits
improved. Every interval includes zero.}
\label{fig:forest}
\end{figure}

\begin{figure}[t]
\centering
\includegraphics[width=\linewidth]{figures/fig2_subgroups_pearson_delta.pdf}
\caption{Pearson~$\Delta$ per test subgroup. No subgroup shows a consistent gain from the learned graph.}
\label{fig:subgroups}
\end{figure}

\begin{figure}[t]
\centering
\includegraphics[width=\linewidth]{figures/fig3_per_perturbation_pearson_delta.pdf}
\caption{Per-perturbation Pearson~$\Delta$, LASSER-fix vs.\ GEARS-fix, pooled over splits. Points scatter around
the diagonal.}
\label{fig:perpert}
\end{figure}

\begin{figure}[t]
\centering
\includegraphics[width=0.85\linewidth]{figures/fig4_gene_groups.pdf}
\caption{Change relative to GEARS-fix on output genes isolated in vs.\ connected in $A_0$. The isolated genes, whose
neighbourhoods change most, do not gain.}
\label{fig:genegroups}
\end{figure}

\begin{figure}[t]
\centering
\includegraphics[width=\linewidth]{figures/fig6_graphs.pdf}
\caption{(a) In-degree distribution on split 1 for the GEARS Pearson graph ($A_0$), the density-matched Pearson
graph and the learned graph. (b) Share of edges with $|r|>0.2$ in held-out cells ($H_2$): $A_0$ edges, learned new
edges and degree-matched random pairs (mean $\pm$ sd over splits).}
\label{fig:graphs}
\end{figure}

\clearpage
\bibliographystyle{unsrtnat}
\bibliography{refs}
\end{document}
